
\section{Basic Trigonometry}
Trigonometry provides tools on the relations between angles and distance ratios in right triangles. 
Given the circle of unit radius ($r=1$) in figure~\ref{fig:trigonometry}, the cosine, sine and tangent are defined as: 
\begin{equation}
 \cos \alpha = \frac{a}{r} = \frac{a}{1} = a; \ \ \ \ 
 \sin \alpha = \frac{b}{r} = \frac{b}{1} = b; \ \ \ \ 
 \tan \alpha = \frac{b}{a};
\end{equation}
The above functions are defined as $f(\alpha):\mathbb{R} \rightarrow \mathbb{R}$. They are periodic functions with period $2\pi$, meaning that the function provides the same value for an argument $\alpha+k2\pi$,~being $k$ an integer ($k \in \mathbb{Z}$). 
\begin{figure}[bth!]
  \begin{center}
    \includegraphics[width=1.0\columnwidth]{figures/trigonometry.eps}
    \caption{Unit circle and deifinitions of $\cos,\ \sin$ and~$\tan$. (to do)}
    \label{fig:trigonometry}
  \end{center}
\end{figure}

The inverse functions $acos(),\ asin()$ and~$atan()$ provides an angle given a ratio of distances. 

% \begin{minipage}[1pt]{c}
%  hello
% \end{minipage} 
